% !TEX root = main.tex
% APPENDIX A (one of several). Use appendixA.tex, appendixB.tex, ... when the
% document has TWO OR MORE appendices (they are lettered A, B, ...).
% Start each file with \chapter{Title}; \appendix in main.tex switches
% chapters to appendix lettering. Every appendix needs a title and a short
% introduction before any figures or tables (GSO).
% The sample text is shared with appendix.tex through appendix-rules.tex and
% appendix-demo.tex.

\chapter{Including an appendix or appendices}
\label{app:appendixA}

% !TEX root = main.tex
% Shared text: the GSO rules for appendices. Included with \input from both
% appendix.tex (single appendix) and appendixA.tex (several appendices), so
% the wording lives in one place. Delete this file and its \input lines
% when you write your own appendices.

An appendix holds material that supports the dissertation but that a reader does not need in order to follow the argument. A useful test: if a reader skipped the appendix, would the chapters still make sense and still support their conclusions? If not, the material belongs in a chapter.

\textbf{Keep in the chapters} everything a reader needs to follow and evaluate the work:
\begin{itemize}
\item the research questions, the methods, and the main results;
\item the figures and tables that present those results;
\item the key steps of the central proofs and derivations.
\end{itemize}

\textbf{Move to an appendix} material that is long, routine, or of interest mainly to a specialist checking the details:
\begin{itemize}
\item long or routine proofs and derivations, once the chapter has stated the result;
\item full data tables, when the chapter shows a summary;
\item additional figures that repeat a result for other parameters or cases;
\item survey instruments, interview protocols, and consent forms;
\item long source code or pseudocode;
\item details of experimental setups, calibrations, and numerical parameters.
\end{itemize}

Refer to every appendix from the chapter where its material is relevant, for example ``The full proof is given in Appendix~B.'' Material that no chapter refers to probably does not belong in the dissertation at all. Large data sets and complete software packages usually do not belong in an appendix either; deposit them in a public repository and cite them, or ask the GSO how to submit supplementary files.

NJIT enforces a strict rule regarding the titling and numbering of appendices:
\begin{itemize}
\item If your dissertation has only a \textbf{single appendix}, it must be titled \textbf{APPENDIX} (without any letter label such as "A"), followed by the appendix title.
\item If your dissertation contains \textbf{two or more appendices}, they must be numbered sequentially as \textbf{APPENDIX A}, \textbf{APPENDIX B}, etc.
\end{itemize}

You should not try to format this manually. Instead:
\begin{itemize}
\item If you have \textbf{one} appendix, use \texttt{appendix.tex} and begin it with the command \verb+\oneappendix{Your Title}+. The heading prints as APPENDIX followed by your title in capitals.
\item If you have \textbf{two or more}, switch to the multi-appendix lines in \texttt{main.tex} and begin each file with \verb+\chapter{Your Title}+. They are lettered A, B, \ldots\ automatically.
\end{itemize}

Every appendix needs a title and a brief introductory description before any equations, figures, or tables (GSO). Appendix pages are numbered continuously with the chapters, and equations, figures, and tables in an appendix are numbered with the appendix letter, starting again at 1 (A.1, B.1, \ldots).


This sample has five appendices, so they are lettered: this is \textbf{APPENDIX A}, followed by \cref{app:citestyles}, which collects guides to the citation styles, \cref{app:packages}, which describes the packages the template loads, \cref{app:filetree}, which lists the template's files, and \cref{app:supplementary}, which lists supplementary files. The class file prints the letters in the headings and the Table of Contents and prefixes equations, theorems, figures, and tables with the letter, as the examples below show. To switch to a single, unlettered appendix, follow the instructions at the end of \texttt{main.tex}.

% !TeX root = main.tex
% Shared sample: demonstrates numbered equations, figures, and a table inside
% an appendix. Included with \input from appendix.tex (single appendix) and
% appendixA.tex (lettered appendices); only one of those is compiled at a
% time, so each label is defined once. Delete this file and its \input lines
% when you write your own appendices.

\section{Equations and Figures in Appendices}

Equations inside the appendix are numbered with an appendix prefix:
\begin{align}
\label{eq:app1}
\nabla^2 \phi - \frac{1}{c^2}\frac{\partial^2 \phi}{\partial t^2} &= 0, \\
\label{eq:app2}
\oint_{\partial \Omega} \mathbf{F} \cdot d\mathbf{S} &= \iiint_{\Omega} (\nabla \cdot \mathbf{F})\, dV.
\end{align}
Cross-references like Equation~\eqref{eq:app1} and Equation~\eqref{eq:app2} resolve cleanly.

Figures placed within the appendix also receive appropriate appendix numbering, as shown in Figure~\ref{fig:app_cube1} and Figure~\ref{fig:app_cube2}.

\begin{figure}[htbp]
\centering
\includegraphics[width=2.5in]{msgcube}
\caption{First appendix demonstration figure.}
\label{fig:app_cube1}
\end{figure}

\begin{figure}[htbp]
\centering
\includegraphics[width=2.5in]{msgcube}
\caption{Second appendix demonstration figure.}
\label{fig:app_cube2}
\end{figure}

\clearpage
\section{Tables in an appendix}

Here is an example of a table in an appendix showing how you might typeset a table of websites you consulted in your thesis. 
Table~\ref{table:WEBSITES} lists common bioinformatics and signaling databases, formatted with automatic line wrapping for URLs using the \texttt{xurl} package.

\begin{table}[htbp]
\centering
\caption{Key Online Databases for Cell Signaling and Molecular Biology}
\label{table:WEBSITES}
\begin{tabularx}{\linewidth}{|c|p{0.32\linewidth}|X|}
\hline
\textbf{\#} & \textbf{Resource / URL} & \textbf{Description} \\
\hline
1 & \url{http://vlib.org/Science/Cell_Biology} & Virtual Library for Cell Biology with curated signal transduction pathways. \\
\hline
2 & \url{http://www.grt.kyushu-u.ac.jp/spad/} & Signaling Pathway Database containing curated pathway diagrams. \\
\hline
3 & \url{http://geo.nihs.go.jp/csndb/} & Cell Signaling Networks Database (CSNDB) \cite{CSND:1999}. \\
\hline
4 & \url{http://www.sdsc.edu/kinases/} & Protein Kinase Resource (PKR) for protein kinase structural data \cite{PKR:1997}. \\
\hline
5 & \url{http://www.expasy.ch/sprot/} & SWISS-PROT annotated protein sequence database \cite{BAIROCH:2000}. \\
\hline
6 & \url{http://www.expasy.ch/prosite/} & PROSITE database of protein families and functional domains \cite{HOFMANN:1999}. \\
\hline
7 & \url{http://www.rcsb.org/pdb/} & Protein Data Bank (PDB) for 3D macromolecular structures \cite{BERMAN:2000}. \\
\hline
\end{tabularx}
\end{table}

